%-------------------------------------------------------------------------------
%	SECTION TITLE
%-------------------------------------------------------------------------------
\cvsection{Experience}

%-------------------------------------------------------------------------------
%	CONTENT
%-------------------------------------------------------------------------------
\begin{cventries}
%---------------------------------------------------------
  \cventry
    {University of Montpellier -- LIRMM -- BIATSS} % Organization
    {Research Engineer} % Job title
    {Montpellier, France} % Location
    {Feb. 2024 -- Feb. 2025} % Date(s)
    {
      \begin{cvitems} % Description(s) of tasks/responsibilities
        \item Maintainer of the control framework \entrypositionstyle{mc\_rtc}
        \item Development and maintenance of LIRMM's robotic platforms, notably humanoid robots (\entrypositionstyle{HRP-4}, \entrypositionstyle{RHPS1}, \entrypositionstyle{Unitree G1}) and robotic arms (\entrypositionstyle{Franka Emika Panda}, \entrypositionstyle{Universal Robots}).
        \item Rollkneematics project: Prototyping a new capacitive sensor for knee prosthesis detection in collaboration with BoneTag.
        \item Development of reproducible demonstrators: dance \entrypositionstyle{HRP-4, RHPS1}, whole-body dynamic stability control of a humanoid, etc.
        \item Support and development for projects of all LIRMM robotics teams, notably humanoid and underwater.
        \item 3D part modeling and printing: sensor integration, mechanical adaptations, etc.
      \end{cvitems}
    }
%---------------------------------------------------------
  \cventry
    {CNRS -- LIRMM -- Interactive Digital Human} % Organization
    {Research Engineer} % Job title
    {Montpellier, France} % Location
    {Nov. 2021 -- April 2023} % Date(s)
    {
      \begin{cvitems} % Description(s) of tasks/responsibilities
        \item Development of the \entrypositionstyle{mc\_rtc} framework
        \item Creation of industrial demonstrators (\emph{confidential})
        \item Participation in the \entrypositionstyle{I.AM} project
        \item Prototyping a new capacitive sensor for knee prosthesis detection in collaboration with BoneTag.
        \item Supervision and support for PhD students and postdocs
        \item Finalist in the international competition \entrypositionstyle{X-Prize} as part of the \entrypositionstyle{Janus} team
      \end{cvitems}
    }

%---------------------------------------------------------
  \cventry
    {Joint Robotics Laboratory -- Advanced Institute of Science of Technology} % Organization
    {Research Engineer} % Job title
    {Tsukuba, Japan} % Location
    {Nov. 2019 -- April 2021} % Date(s)
    {
      \begin{cvitems} % Description(s) of tasks/responsibilities
        \item Responsible for the unification and merging of control software between~:\\
          -- The \entrypositionstyle{mc\_rtc} framework developed by CNRS by the IDH team at LIRMM (Montpellier) and AIST-JRL (Tsukuba)\\
          -- The \entrypositionstyle{HMC} framework developed by the HRG group at AIST (Tsukuba)
        \item Technical support for both groups, development of demonstrations on robots to ensure experimental contributions to scientific publications as well as to meet the expectations expressed by the many bilateral projects with industry partners.
      \end{cvitems}
    }

%---------------------------------------------------------
  \cventry
    {CNRS -- LIRMM -- Interactive Digital Human} % Organization
    {Research Engineer} % Job title
    {Montpellier, France} % Location
    {Oct. 2018 - Oct. 2019} % Date(s)
    {
      \begin{cvitems} % Description(s) of tasks/responsibilities
        \item \entrypositionstyle{H2020 COMANOID - Multi-Contact Collaborative Humanoids in Aircraft Manufacturing}\\
          \entrypositionstyle{Site}: \url{https://comanoid.cnrs.fr/}\\
          \entrypositionstyle{Role}: \textbf{Responsible for the implementation and integration of localization and mapping methods} for the final demonstration of the European COMANOID project. This demonstration, the result of 4 years of joint effort between four research institutes (LIRMM, DLR, University of Rome La Sapienza, Inria Rennes and Grenoble) demonstrated the applicability of manufacturing humanoid robots in the real industrial context of aircraft construction. This mainly involved locomotion and manipulation capabilities in a constrained space: walking and localization (SLAM), stair climbing (MPC), manipulation (SLAM, registration, visual servoing, force control, etc.).
        \item \entrypositionstyle{mc\_rtc}: \url{https://jrl-umi3218.github.io/mc_rtc}\\
        Development and maintenance of the control framework {\tt mc\_rtc} used in the above demonstrator, as well as by students and researchers at LIRMM, JRL, and their partners.
        \item Technical support for IDH students and researchers and conducting experiments on HRP-4 and BAZAR robots.
      \end{cvitems}
    }

%---------------------------------------------------------
  \cventry
    {LIRMM, I3S, JRL} % Organization
    {PhD} % Job title
    {France, Japan} % Location
    {Oct. 2014 - Nov. 2018} % Date(s)
    {
      \begin{cvitems} % Description(s) of tasks/responsibilities
        \item \entrypositionstyle{Supervisors}: Abderrahmane Kheddar, Andrew Comport 
        \item \entrypositionstyle{Projects}: RobotHow, H2020 COMANOID, DARPA Robotics Challenge
        \item Localization of a humanoid robot and its environment using state-of-the-art dense visual SLAM.
        \item Object localization by registering CAD models with the dense SLAM map.
        \item Online adaptation of offline-generated multi-contact locomotion plans using SLAM-based localization and mapping.
        \item Development of a whole-body calibration method.
        \item Walking by model predictive control (MPC), using a fusion of visual (SLAM) and proprioceptive (encoders, force sensors) information to react to disturbances by continuously generating a ZMP trajectory and future steps ensuring the robot's stability.
        \item \entrypositionstyle{\emph{DARPA Robotics Challenge (DRC)}}: Participation as part of the AIST-NEDO team. Ranked 10/23 with completion of 6 out of 8 tasks (semi-autonomous driving, opening a door and a valve, drilling a wall, connecting a cable, traversing rough terrain).
        \item \entrypositionstyle{Software contributions}\\
          - Registration methods used during the DARPA Robotics Challenge: 3D point cloud registration (\url{https://github.com/arntanguy/icp}), generation of 3D point clouds from their CAD models (\url{https://github.com/arntanguy/mesh\_sampling})\\
          - Savitzky-Golay filter -- \url{https://github.com/arntanguy/gram_savitzky_golay}\\
          - Whole-body calibration -- \url{https://github.com/arntanguy/robcalib}\\
          - Contributions to {\tt mc\_rtc}: state observation (SLAM, IMU, VICON ground truth,...), trajectory tracking tasks, etc.
      \end{cvitems}
    }

%---------------------------------------------------------
  \cventry
    {Technische Universität München (TUM)} % Organization
    {Intern} % Job title
    {Munich, Germany} % Location
    {2014 (6 months)} % Date(s)
    {
      \begin{cvitems} % Description(s) of tasks/responsibilities
        \item \entrypositionstyle{Supervisors}: Jurgen Sturm and Daniel Cremers
        \item Application of convolutional neural networks to loop closure detection in SLAM.
        \item Development of the architecture enabling the use of Siamese networks in the open-source framework {\tt Caffe}.
      \end{cvitems}
    }

%---------------------------------------------------------
  \cventry
    {Polytech Nice-Sophia-Antipolis, Trinity College Dublin} % Organization
    {University projects} % Job title
    {France, Ireland, Germany} % Location
    {2014 (6 months)} % Date(s)
    {
      \begin{cvitems} % Description(s) of tasks/responsibilities
        \item Development of a physics and rendering engine (fluid simulation, rigid body collisions, object/fluid collisions, raytracing)\\\url{https://github.com/arntanguy/PHEngine}.
        \item Development of an interactive curve fitting software specialized for research in scanning tunneling microscopy spectroscopy\\\url{https://github.com/arntanguy/STS-simulator}. 
        \item Photorealistic rendering of SLAM maps in an Occulus Rift (project led by Andrew Comport).
        \item Development of a 3D racing game for visually impaired players\\\url{http://prdevint.polytech.unice.fr}. 
        \item Development of augmented reality games.
      \end{cvitems}
    }

%---------------------------------------------------------
  \cventry
    {Fotowall} % Organization
    {High school student, self-taught C++ project} % Job title
    {Brest, France} % Location
    {2008-2011} % Date(s)
    {
      \begin{cvitems} % Description(s) of tasks/responsibilities
        \item \entrypositionstyle{Site}: \url{https://www.enricoros.com/opensource/fotowall/index.html}
        \item Self-taught development of the open-source C++ image manipulation software Fotowall
        \item Remote collaboration with the Italian developer Enrico Ross
        \item Over one million downloads (as of 2017)
      \end{cvitems}
    }

\end{cventries}
